\section{Category Methods and Data Requirements}
\label{sec:factors}

The framework defines an IPCC calculation method and an evidence requirement for each
category. Higher-resolution activity data can therefore replace aggregate statistics without
changing the material balance or category definition.

\textbf{Mineral industries.} Cement uses Equation~\ref{eq:clinkerbalance}. Generic carbonate
methods are specified for lime, glass and other carbonate uses, but their Pakistan rows remain
NE until lime output by type, glass mass and type, cullet use and carbonate inputs are compiled.

\textbf{Chemical industries.} The ammonia ledger divides gross process \COtwo{} among
atmospheric emissions, recovery into urea, other recovery and permanent storage and asserts
zero residual. Nitric-acid emissions are estimated from production expressed as 100\% acid,
technology-specific
\NtwoO{} factors, the abated share, destruction efficiency and abatement utilisation. Total
fertiliser output or nameplate capacity alone is insufficient for either estimate.

\textbf{Metals and non-energy products.} QIM reports billet, ingot and lubricant production.
Iron and steel estimation additionally requires the production route and carbon inputs;
lubricant emissions require non-energy consumption and oxidation. These remaining inputs are
documented in the completeness matrix.

\textbf{Fluorinated gases.} A species-resolved equipment-bank model conserves annual chemical
supply across manufacturing loss, installed bank, operating leakage, retirement, recovery,
destruction and recovered gas available for reuse. AR5 and AR6 GWP registries include the HFC,
PFC, SF$_6$ and NF$_3$ species required by the CRT. The Ozone profile lacks the species- and
application-level data needed to estimate the bank, so 2.F remains NE in the independent
FY2023--24 inventory.
